\documentclass[a4paper,12pt]{exam}
\usepackage[utf8]{inputenc}


\begin{document}
	
\begin{center}
	\fbox{\fbox{\parbox{5.5in}{\centering
				This is a multiple choice practice exam. Time yourself for 20 minutes and answer the following questions.}}}
	\\
	\vspace{5mm}
	{\LARGE \textbf{Day 1 Biology Homework}}
\end{center}

\vspace{1mm}
\makebox[\textwidth]{Name:\enspace\hrulefill}

\vspace{5mm}

\makebox[\textwidth]{Date:\enspace\hrulefill}

\begin{questions}
	\question What is the first step in identifying the scientific method?

	\begin{oneparchoices}
		\choice Identify the question
		\\
		\choice Formulate a hypothesis
		\\
		\choice Collect and analyze data
		\\		
		\choice Make conclusions
	\end{oneparchoices}

	\question Edward wants to find the effect of saltwater on plants, so he waters one plant with saltwater each day for a month, and does not water another plant for a month. Both plants are kept in the same environmental conditions. What is the dependent variable in this experiment?
	
	\begin{oneparchoices}
		\choice Whether Edward waters the plant with saltwater
		\\
		\choice The amount of water that is given to the plant with saltwater
		\\
		\choice The amount of sunlight each plant is given
		\\		
		\choice The growth of the plants
	\end{oneparchoices}

	\question Who developed the system of binomial nomenclature, a two-word system for naming all species on Earth?
	
	\begin{oneparchoices}
		\choice James Watson
		\\
		\choice Carolus Linneas
		\\
		\choice Antonie van Leeuwenhoek
		\\
		\choice Marie Curie
	\end{oneparchoices}

	\question What is the format for binomial nomenclature, in which "genus" refers to the genus of the organism, and "species" refers to the species of the organism?
	
	\begin{oneparchoices}
		\choice \textit{species genus}
		\\
		\choice \textit{Genus species}
		\\
		\choice \textit{Species genus}
		\\
		\choice \textit{genus species}
	\end{oneparchoices}

	\question What is the process by which larger molecules are broken down into smaller molecules with the addition of water?
	
	\begin{oneparchoices}
		\choice Dehydration synthesis
		\\
		\choice Chemical digestion
		\\
		\choice Hydrolyis
		\\
		\choice Mechanical digestion
	\end{oneparchoices}

	\question What is the name of the smallest vessels in the circulatory system?
	
	\begin{oneparchoices}
		\choice Veins
		\\
		\choice Aorta
		\\
		\choice Arterioles
		\\
		\choice Capillaries
	\end{oneparchoices}

	\question What is the name of the disks that contain the pigment chlorophyll in the chloroplast?
	
	\begin{oneparchoices}
		\choice Granum
		\\
		\choice Thylakoids
		\\
		\choice Stroma
		\\
		\choice Inner Membrane
	\end{oneparchoices}

	\question Determine the steps of photosynthesis in order.
	
	\begin{oneparchoices}
		\choice Light-Dependent Reactions, Calvin Cycle
		\\
		\choice Light-Dependent Reactions, Glycolysis, Calvin Cycle
		\\
		\choice Calvin Cycle, Light-Dependent Reactions
		\\
		\choice Krebs Cycle, Oxidative Phosphorylation, Light-Dependent Reactions
	\end{oneparchoices}

	\question Where is the electron transport chain located?
	
	\begin{oneparchoices}
		\choice Intermembrane space in the mitrochondria
		\\
		\choice Inner mitochondrial membrane
		\\
		\choice Cytosol
		\\
		\choice Nucleus
	\end{oneparchoices}

	\question What are the 2 types of fermentation that occur in anaerobic respiration?
	
	\begin{oneparchoices}
		\choice Lactic Acid Fermentation and Krebs Cycle
		\\
		\choice Lactic Acid Fermentation and Alcohol Fermentation
		\\
		\choice Pyruvate and Glycolysis
		\\
		\choice Oxaloacetatic Fermentation and Citrate Cycle
	\end{oneparchoices}
	
	\question What is one function of lipids?
	
	\begin{oneparchoices}
		\choice Long-term energy storage
		\\
		\choice Phospholipids in the cell membrane
		\\
		\choice Steroids
		\\
		\choice All of the above choices are functions of lipids
	\end{oneparchoices}

	\question What are the components of nucleic acids?
	
	\begin{oneparchoices}
		\choice Phosphate group, pentose sugar, nitrogenous base
		\\
		\choice Amino group, carboxyl group, R group
		\\
		\choice Glycerol, fatty acids
		\\
		\choice Glucose, fructose, maltose, sucrose
	\end{oneparchoices}

	\question Name the process which is responsible for the replication of sex/gamete cells.
	
	\begin{oneparchoices}
		\choice Mitosis
		\\
		\choice Meiosis
		\\
		\choice Cytokinesis
		\\
		\choice Metaphase
	\end{oneparchoices}

	\question Define codominance.
	
	\begin{oneparchoices}
		\choice The dominant and recessive traits of a heterozygous trait gets blended together
		\\
		\choice There are more than 2 genotypes for an organism
		\\
		\choice The dominant and recessive traits of a heterozygous trait both are expressed separately
		\\
		\choice Traits that will be only be expressed when an organism does not have a conflicting dominant trait
	\end{oneparchoices}

	\question A fire occurs in a forest, burning everything to the ground. What is the first thing that will happen as the forest attempts to regrow?
	
	\begin{oneparchoices}
		\choice Bare rock is broken down by lichens to create soil
		\\
		\choice Large trees immediately begin to grow
		\\
		\choice Small grasses and perennials begin to appear
		\\
		\choice Pioneer species start to grow
	\end{oneparchoices}

	\question A small, charged molecule attempts to pass through the cell membrane. Determine what happens to the molecule immediately.
	
	\begin{oneparchoices}
		\choice The molecule is not allowed to pass through the cell membrane
		\\
		\choice The molecule is stored in the cell membrane 
		\\
		\choice The molecule freely passes through the cell membrane
		\\
		\choice The charge of the molecule is removed and the molecule passes through the membrane 
	\end{oneparchoices}

	\question What is an example of passive transport?
	
	\begin{oneparchoices}
		\choice Osmosis
		\\
		\choice Sodium-Potassium Pump
		\\
		\choice Proton Pump
		\\
		\choice Endocytosis
	\end{oneparchoices}

	\question What is the process by which a cell ingests large particles?
	
	\begin{oneparchoices}
		\choice Pinocytosis
		\\
		\choice Phagocytosis
		\\
		\choice Exocytosis
		\\
		\choice Receptor-mediated endocytosis
	\end{oneparchoices}

	\question How many bones does the appendicular skeleton contain?
	
	\begin{oneparchoices}
		\choice 126 bones
		\\
		\choice 80 bones
		\\
		\choice 206 bones
		\\
		\choice 412 bones
	\end{oneparchoices}

	\question A fire occurs in a forest, burning everything to the ground. What is the first thing that will happen as the forest attempts to regrow?
	
	\begin{oneparchoices}
		\choice Bare rock is broken down by lichens to create soil
		\\
		\choice Large trees immediately begin to grow
		\\
		\choice Small grasses and perennials begin to appear
		\\
		\choice Pioneer species start to grow
	\end{oneparchoices}

	\question What is the name of the muscle that is found in the small intestine?
	
	\begin{oneparchoices}
		\choice Skeletal muscle
		\\
		\choice Appendicular muscle
		\\
		\choice Smooth Muscle
		\\
		\choice Cardiac Muscle
	\end{oneparchoices}

	\question Name a function of the integumentary system.
	
	\begin{oneparchoices}
		\choice Secretes hormones to maintain homeostasis
		\\
		\choice Detects pathogens that may harm the body
		\\
		\choice Collects excess fluid and returns it to the blood
		\\
		\choice Regulation of body temperature
	\end{oneparchoices}

	\question What are cell walls composed of?
	
	\begin{oneparchoices}
		\choice Starch
		\\
		\choice Glucose
		\\
		\choice Cellulose
		\\
		\choice Maltose
	\end{oneparchoices}

	\question Fill in the blanks: \_\_\_\_\_\_\_\_  carries water and minerals upwards from the roots, and \_\_\_\_\_\_\_ carries sugars made by photosynthesis to leaves.
	
	\begin{oneparchoices}
		\choice Xylem, phloem
		\\
		\choice Xylem, vacuole
		\\
		\choice Vascular system, non-vascular system
		\\
		\choice Spores, rhizome
	\end{oneparchoices}

	
	\question Name an example of an angiosperm.
	
	\begin{oneparchoices}
		\choice Sequoia
		\\
		\choice Redwood
		\\
		\choice Mosses
		\\
		\choice Dandelions
	\end{oneparchoices}

	
	\question What are the products of cellular respiration?
	
	\begin{oneparchoices}
		\choice Carbon and water
		\\
		\choice Glucose and oxygen
		\\
		\choice ATP, water, and carbon dioxide
		\\
		\choice Chlorophyll, oxygen, and chitin
	\end{oneparchoices}

	
	\question Which of the following has only the elements carbon, hydrogen, and oxygen in the ratio 1:2:1, respectively?
	
	\begin{oneparchoices}
		\choice Proteins
		\\
		\choice Lipids
		\\
		\choice Carbohydrates
		\\
		\choice Nucleic Acids
	\end{oneparchoices}

	
	\question Who is known as the "father of genetics"?
	
	\begin{oneparchoices}
		\choice Richard Feynman
		\\
		\choice Charles Darwin
		\\
		\choice Rosalind Franklin
		\\
		\choice Gregor Mendel
	\end{oneparchoices}

	
	\question When is a solution hypertonic with respect to a cell placed in it?
	
	\begin{oneparchoices}
		\choice The solution has a higher solute concentration than the cell
		\\
		\choice The concentration of solute in the cell and in the solution is the same
		\\
		\choice The cell has a higher solute concentration than the solution
		\\
		\choice The solution has a lower solute concentration than the cell
	\end{oneparchoices}

	
	\question What begins a feedback mechanism?
	
	\begin{oneparchoices}
		\choice Homeostatic disruption
		\\
		\choice Homeostasis monitor
		\\
		\choice External stimulus
		\\
		\choice Receptor (sensor)
	\end{oneparchoices}
\end{questions}

\end{document}